\section{A general period and all special fibers}
\label{spread:section}

The construction has produced an algebraic class outside the span of
polarization powers whenever the generic theta algebra is realized at a
suitable complex period.  Two steps remain.  We choose such a period
where this detects a nonzero Weil component, and then use a projective
family of subschemes to carry one fixed Weil class to every period.
The exceptional sets used for this choice must be fixed before the
matrix direction $M$ is selected.

\subsection{Countably many projective parameter spaces}

Let $\pi:\mathscr A\to\mathcal U$ be the holomorphic family in
\eqref{hodge:period-domain}.  Its underlying topological family is the
marked product $T\times\mathcal U$.  The $h_0$th power of the theta
line has a holomorphic basis of theta sections over $\mathcal U$;
their multiplier formulas and normal convergence are those used in
Section~\ref{theta:section}.  Since $h_0\ge3$, this power is very ample
on every abelian fiber \cite{abelian-ampleness}.  These sections give a
closed holomorphic embedding
\begin{equation}
 \mathscr A\hookrightarrow\mathbb P^N\times\mathcal U,
 \qquad N=h_0^8-1.
 \label{spread:relative-embedding}
\end{equation}
Indeed the map is proper and fiberwise an embedding; its differential
is injective on vertical tangents, and its base map is the identity.

\begin{lemma}[Hilbert parameter spaces]
\label{spread:hilbert-loci}
There is a countable collection of connected reduced complex analytic
spaces $\mathcal H$, each proper over $\mathcal U$, with a flat family
of subschemes
\[
 \mathscr Z_{\mathcal H}\subset
       \mathscr A\times_{\mathcal U}\mathcal H,
\]
such that every closed subscheme of every fiber $A_\Pi$ occurs in at
least one family.  Each image $\mathcal H\to\mathcal U$ is closed
analytic.  Consequently the union $\mathcal E_{\mathrm{Hilb}}$ of
those images which are proper subsets is a countable union of proper
closed analytic subsets of $\mathcal U$.
\end{lemma}

\begin{proof}
For each Hilbert polynomial $P$, take the projective Hilbert scheme
$\operatorname{Hilb}_P(\mathbb P^N)$ and its universal flat family
\cite{grothendieck-hilbert}.  Write $H_P$ for its analytification.  We
first construct the reduced closed analytic incidence locus
\[
 \mathcal H_P=\{(\Pi,z)\in\mathcal U\times H_P:
                     Z_z\subset A_\Pi\text{ scheme-theoretically}\}.
\]
Here the scheme structure of $Z_z$ is retained.

Fix a small neighborhood $V$ of a period.  The ideal of $\mathscr A$
in $\mathbb P^N\times V$ is flat over $V$, since both the ambient
structure sheaf and its quotient are flat.  For a sufficiently large
twist $a$, relative generation and proper analytic cohomology and base
change give finitely many homogeneous polynomial families
$s_1,\ldots,s_r$ of degree $a$ generating that ideal after shrinking
$V$ \cite{grauert-remmert}.  One can see the last assertion directly
from the evaluation map: it is surjective on the central projective
fiber, and the support of its coherent cokernel has proper closed image,
which is removed by shrinking.

On the fixed projective Hilbert scheme choose $e\ge0$ large enough that
the universal family's sheaf of degree-$a+e$ sections pushes forward to
a vector bundle $E_P$ and commutes with base change.  One $e$ works
over the entire Hilbert scheme by relative Serre vanishing and its
projectivity.  Restrict each product $s_i x^\beta$, with $|\beta|=e$,
to the pulled-back universal family over $V\times H_P$.  These
restrictions are holomorphic sections of the pulled-back bundle $E_P$.
Their common zero locus has exactly the required incidence points.
For the nontrivial implication, on the projective chart $x_j\ne0$ the
tested product $s_i x_j^e$ vanishes, so $s_i$ vanishes in the structure
sheaf of the subscheme.  The $s_i$ generate the ideal of the abelian
family, which proves scheme-theoretic containment.

This also proves containment over the entire zero locus: base change
identifies the zero testing sections with their evaluations on its
universal family, and the same chart argument annihilates the pulled-back
ideal.  Taking reduced structures, the local zero loci agree on overlaps
and give $\mathcal H_P$.  Pullback of the universal family is still flat
and is a subscheme of the pulled-back abelian family.  Reduction of the
parameter space does not require any fiber subscheme to be reduced.

The projection $\mathcal H_P\to\mathcal U$ is proper because $H_P$ is
compact.  Complex analytic spaces are locally path connected, so each
connected component of $\mathcal H_P$ is open and closed, in particular
a closed analytic subspace.  There are at most countably many components:
the space is second countable, and each nonempty open component contains
a member of a countable basis.  The components remain proper over
$\mathcal U$, and Remmert's proper mapping theorem makes their images
closed analytic \cite{grauert-remmert}.  Finally, there are countably
many Hilbert polynomials.  These components are the asserted spaces.
\end{proof}

\subsection{Constancy across singular parameters}

The parameter spaces just constructed need not be smooth.  We prove
constancy of their cohomology classes in a form which includes such
parameters and nonreduced specialized subschemes.

\begin{lemma}
\label{spread:chern-constancy}
For every connected parameter space $\mathcal H$ in
Lemma~\ref{spread:hilbert-loci}, the marked rational cohomology class
\[
 \ch(\mathcal O_{Z_b})\in H^{\mathrm{even}}(T,\Q)
\]
is independent of $b\in\mathcal H$.  In particular its degree-eight
component $\ch_4(\mathcal O_{Z_b})$ is constant.  On each fiber it is
a rational algebraic cycle class.
\end{lemma}

\begin{proof}
Put $X=\mathscr A\times_{\mathcal U}\mathcal H$ and
$\mathcal F_0=\mathcal O_{\mathscr Z_{\mathcal H}}$, considered as a
coherent sheaf on $X$.  It is flat over $\mathcal H$.  Near any chosen
$b\in\mathcal H$, construct eight exact sequences
\begin{equation}
 0\longrightarrow\mathcal F_{j+1}\longrightarrow E_j
   \longrightarrow\mathcal F_j\longrightarrow0,
 \qquad 0\le j\le7,
 \label{spread:syzygies}
\end{equation}
where $E_j$ is a vector bundle and every $\mathcal F_j$ is flat over
the parameter neighborhood.  To obtain a sequence, choose a sufficiently
positive twist of $\mathcal F_j$ on the central projective fiber so that
it is generated by its sections and has no higher cohomology.  Proper
analytic cohomology and base change give a vector bundle of sections
on a smaller neighborhood.  Its evaluation, tensored by the opposite
twist, is a surjection near the whole central fiber, and hence over
a smaller base neighborhood by properness.  Since $X$ is flat over
the base, $E_j$ is base-flat; the kernel in \eqref{spread:syzygies}
is base-flat as well.  In particular these sequences restrict exactly
to each fiber.

The smooth fiber $X_b$ has dimension eight.  Its local rings are regular
of dimension at most eight, so the eighth syzygy
$\mathcal F_8|_{X_b}$ is locally free.  This property lifts to a
neighborhood of the fiber even though $\mathcal H$ may be singular.
Here is the local argument.  At a point $x\in X_b$ put
\[
 A=\mathcal O_{\mathcal H,b},\quad
 S=\mathcal O_{X,x},\quad M=(\mathcal F_8)_x.
\]
The module $M$ is finite over the Noetherian local ring $S$, is flat
over $A$, and $M/\mathfrak m_A M$ is free over
$S/\mathfrak m_A S$.  Lift a fiber basis to a map $S^r\to M$.
Nakayama makes it surjective.  If $K'$ is its kernel, flatness of $M$
over $A$ makes reduction exact on the left.  The reduced map is an
isomorphism, hence $K'/\mathfrak m_A K'=0$, and another application
of Nakayama gives $K'=0$ \cite{stacks-flat-free}.  Thus $M$ is free.

Local freeness is open.  Properness of $X\to\mathcal H$ now allows one
more shrink of the parameter neighborhood, removing the image of the
nonfree locus of $\mathcal F_8$.  We have obtained a finite vector-bundle
resolution on the whole inverse image of that neighborhood:
\begin{equation}
 0\longrightarrow\mathcal F_8\longrightarrow E_7\longrightarrow
 \cdots\longrightarrow E_0\longrightarrow\mathcal F_0
 \longrightarrow0.
 \label{spread:relative-resolution}
\end{equation}
Only the fibers were required to be regular; no regularity of the total
space over a singular parameter was used.

All the bundles in \eqref{spread:relative-resolution} are complex
topological vector bundles on the marked product
$T\times|\mathcal H|$.  Along a path in the chosen parameter neighborhood,
their restrictions to its two endpoints have equal Chern characters
in $H^*(T,\Q)$, by homotopy invariance.  Exact restriction of
\eqref{spread:relative-resolution} identifies the alternating sum with
$\ch(\mathcal O_{Z_b})$ on each fiber.  The class is therefore locally
constant.  Local path connectedness and connectedness of $\mathcal H$
make it constant on the whole component.  On each smooth projective
fiber the same resolution computes the usual algebraic Chern character
of the coherent structure sheaf, whose cohomology class is the one just
computed.
\end{proof}

\subsection{Choosing the matrix direction}

Let $\mathcal E_{\mathrm{Hdg}}$ be the exceptional union in
Proposition~\ref{hodge:generic-hodge}, and set
\[
 \mathcal E=\mathcal E_{\mathrm{Hdg}}\cup
               \mathcal E_{\mathrm{Hilb}}.
\]
Each constituent is a proper closed analytic subset of the connected
complex manifold $\mathcal U$, hence has empty interior.  Baire's
theorem supplies a point $\Pi_*\in\mathcal U\setminus\mathcal E$.
Choose, before performing the construction,
\begin{equation}
 M=2\pi i(\Pi_*-i1_8).
 \label{spread:direction}
\end{equation}
It is symmetric and commutes with $D$, exactly as required for the
relative two-form $B=dq^tM\wedge dp$.

\begin{proposition}
\label{spread:period-selection}
For this choice of $M$, there is a period $\Pi\notin\mathcal E$
and a codimension-four integral subvariety $Z\subset A_\Pi$ such that
\begin{equation}
 [Z]=a\theta^4+w,\qquad a\in\Q,\quad 0\ne w\in W_K.
 \label{spread:selected-cycle}
\end{equation}
\end{proposition}

\begin{proof}
Use the convex two-parameter domain $\mathcal P_M$ of
\eqref{theta:period-domain}, with period map
\[
 (\tau,s)\longmapsto\tau1_8+sM/(2\pi i).
\]
It sends $(i,1)$ to $\Pi_*$.  Hence the inverse image of every
constituent of $\mathcal E$ is a proper closed analytic subset of
$\mathcal P_M$.  Run the symplectic and algebraic construction for
this fixed $M$.  Proposition~\ref{theta:realization} adds only countably
many proper closed analytic exceptional sets for the prescribed
theta-coefficient embedding.  By Baire, a point avoids both lists.

At its period $\Pi$, Proposition~\ref{lift:complex-class} gives a
rational algebraic total cycle class outside the span of the powers
of $\theta$.  Formula~\eqref{hodge:generic-classes} implies that its
codimension-four component has nonzero projection to $W_K$.
Express a cycle representing this component as a finite rational
linear combination of integral codimension-four subvarieties.  At least
one constituent has a nonzero Weil projection.  Every such constituent
has rational Hodge class and the period lies outside
$\mathcal E_{\mathrm{Hdg}}$, so its class has precisely the form
\eqref{spread:selected-cycle}.
\end{proof}

\subsection{The whole Weil plane on every fiber}

\begin{theorem}
\label{spread:all-fibers}
For every $\Pi\in\mathcal U$, all classes in $W_K(A_\Pi)$ are
rational algebraic.  The same holds for every split Weil-type
polarized abelian eightfold over $\C$, for every imaginary quadratic
field and every compatible polarization type.
\end{theorem}

\begin{proof}
Let $Z$ and $a\theta^4+w$ be as in
Proposition~\ref{spread:period-selection}.  Its Hilbert point belongs
to one of the connected parameter spaces $\mathcal H$ in
Lemma~\ref{spread:hilbert-loci}.  Since $\Pi$ was chosen outside
$\mathcal E_{\mathrm{Hilb}}$, that component's image cannot be a
proper subset of $\mathcal U$.  Thus $\mathcal H\to\mathcal U$
is surjective.

For a reduced integral codimension-four subvariety of a smooth variety,
the leading Chern-character formula is
\[
 \ch_4(\mathcal O_Z)=[Z]
\]
\cite{fulton-intersection}.  Lemma~\ref{spread:chern-constancy}
therefore gives the fixed marked class $a\theta^4+w$ as
$\ch_4(\mathcal O_{Z_b})$ at every point $b\in\mathcal H$.
Over an arbitrary period, choose any point of the surjective parameter
space.  Its subscheme may be singular, reducible, or nonreduced, but
this Chern character is still a rational algebraic cycle class on the
smooth abelian fiber.  Subtracting the rational algebraic class
$a\theta^4$ proves that the same nonzero rational vector $w\in W_K$
is algebraic on every fiber.

To obtain a second independent class, choose an integer
\[
 m>\sqrt d\,\cot(\pi/8).
\]
The integral matrix $m1_8+D$ commutes with every period in $\mathcal U$,
and hence defines an algebraic endomorphism of every $A_\Pi$.  On the
two determinant lines of $W_K\otimes\C$, its pullback has eigenvalues
\[
 \lambda_+=(m+i\sqrt d)^8,\qquad
 \lambda_-=(m-i\sqrt d)^8.
\]
They are nonreal conjugates: the argument of $m+i\sqrt d$ lies
strictly between $0$ and $\pi/8$.  A nonzero rational $w$ has nonzero
components on both conjugate determinant lines.  If its pullback were
a rational multiple $r w$, those two components would give
$\lambda_+=r=\lambda_-$, a contradiction.  Thus $w$ and its pullback
are independent rational algebraic classes and span $W_K$.

The integer $d$ was arbitrary.  Finally
Proposition~\ref{hodge:isogeny-reduction} transports the conclusion
to every tuple in the theorem, including every polarization type and
every special member with additional endomorphisms.
\end{proof}
